\documentclass[book.tex]{subfiles}
\begin{document}

\chapter{Tensor Field Networks}
\label{chap:neural_equivariant_potentials}
\noindent\rule{11cm}{0.4pt} \\
\textbf{Prerequisites:}  \autoref{chap:molecular_hamiltonian},   \\
\textbf{Difficulty Level:} ** \\
\noindent\rule{11cm}{0.4pt}
\vspace{5mm}

The molecular representations we have learned thus far do not respect known physical constraints about molecular systems. In particular, learned representations for a molecule should update correctly under three dimensional rotations, translations and permutations. That is, representations should update in an equivariant fashion. Respecting equivariance natively in an architecture lowers the requirement for data augmentation techniques which manually add rotated datapoints to training data but at the cost of higher learning complexity. 

Tensor field networks \cite{thomas2018tensor} use spherical harmonics to encode rotational equivariance correctly. Tensor field networks introduce some powerful equivariant primitives that are used in other architectures like neural equivariant potentials. In particular, tensor field networks track scalar, vector, and higher order tensors for every atom in the system.


\section{A Brief Review of Equivariance}

A representation $\rho$ for a group $G$ is a function $\rho: G \to \mathbb{R}^{N \times N}$ such that
\begin{align}
    \rho(g)\rho(h) &= \rho(gh)
\end{align}
Suppose that $X$ and $Y$ are vector spaces with representations $\rho^X$, $\rho^Y$. Then a function $L: X \to Y$ is equivariant with respect to group $\rho$ and representations $\rho^X$ and $\rho^Y$ if the following operators are equal
\begin{align}
L \circ \rho^X(g) &= \rho^Y(g) \circ L
\end{align}
A function $L$ is invariant if $\rho^Y(g) = I$ for all $g \in G$. In the exercises, we ask you to prove some composition properties of equivariant networks.
%"""
%We introduce tensor field neural networks, which are locally equivariant to 3D rotations, translations, and permutations of points at every layer. 3D rotation equivariance removes the need for data augmentation to identify features in arbitrary orientations. Our network uses filters built from spherical harmonics; due to the mathematical consequences of this filter choice, each layer accepts as input (and guarantees as output) scalars, vectors, and higher-order tensors, in the geometric sense of these terms. We demonstrate the capabilities of tensor field networks with tasks in geometry, physics, and chemistry.
%"""

\section{Tensor Field Networks}

Each tensor field network layer takes in a set of vectors $S = \{\vec{r}_a\}$ in $\mathbb{R}^3$, along with a vector $x_a \in X$ for each $\vec{r}_a$ in $S$ and outputs a vector $y_a \in Y$. 
\begin{align}
    L(\vec{r}_a, x_a) &= (\vec{r}_a, y_a)
\end{align}
Here the $x_a, y_a$ should be interpreted as feature vectors, and $X$ and $Y$ are spaces of appropriate features.

\subsection{Permutation Equivariance}

We want our layer $L$ to be equivariant to the action of permutations on the atoms. Formally, we can define permutation

\begin{align}
    P_\sigma(\vec{r}_a, x_a) &= (\vec{r}_{\sigma(a)}, x_{\sigma(a)})
\end{align}

Here $P_\sigma$ permutes the atoms, where 
\begin{align}
\sigma: \{1,\dotsc, n\} \to \{1,\dotsc, n\}
\end{align}
is a bijective map. For $L$ to be equivariant to permutations, we must have that
\begin{align}
    P_{\sigma} \circ L &= L \circ P_{\sigma}
\end{align}

\subsection{Translation Equivariance}

We define translations by operators that shift the position of all atoms without changing their feature vectors
\begin{align}
    T_{\vec{t}}(\vec{r}_a, x_a) &= (\vec{r}_a + \vec{t}, x_a)
\end{align}
For $L$ to be equivariant to translations, we impose the condition that
\begin{align}
    T_{\vec{t}} \circ L &= L \circ T_{\vec{t}}
\end{align}

\subsection{Rotation Equivariance}
Recall that $SO(3)$ is the group of three dimensional rotations. Note that $g \in SO(3)$ can act on vector $\vec{r} \in \mathbb{R}^3$ by rotating this vector. We write this action as
\begin{align}
    R(g)\vec{r} \in \mathbb{R}^3
\end{align}
where $R$ is the representation of $SO(3)$ on $\mathbb{R}^{3 \times 3}$. Suppose that we have representation $\rho^X$ of $SO(3)$ on vector space $X$ and $\rho^Y$ on vector space $Y$. Then note that $g \in SO(3)$ can act on a feature vector $x_a$ as
\begin{align}
    \rho^X(g)x_a \in X
\end{align}
We define a joint operator
\begin{align}
    [R(g) \oplus \rho^X(g)](\vec{r}_a, x_a) &= (R(g)\vec{r}_a, \rho^X(g)x_a)
\end{align}
The rotational equivariance condition is then given by
\begin{align}
    L \circ [R(g) \oplus \rho^X(g)] &= [R(g) \oplus \rho^Y(g)] \circ L
\end{align}
We will make use of a particular family of representations, the Wigner $D$-matrices, 
\begin{align}
    D^{(l)}: SO(3) \to \mathbb{R}^{(2l+1)\times(2l+1)}
\end{align}
which maps elements of $SO(3)$ to square matrices of side $2l+1$. We define
\begin{align}
    D^{(0)} &= I \\
    D^{(1)} &= R(g)
\end{align}
The number $l$ is referred to as the rotation order of the representation.

\section{Constructing Tensor Field Layers}

As mentioned earlier, each tensor field network layer takes in a collection $S = \{\vec{r}_a\}$ of points in $\mathbb{R}^3$ and a collection of associated feature vectors $\{x_a\}$ that lives in a vector space $X$ with a representation of $SO(3)$. This representation can be uniquely decomposed as a sum of irreducible representations of $SO(3)$ on this space. We may have multiple copies of a given irreducible representation (say of order $l$) in this decomposition. These copies are called "channels", in analogy with channels in convolutional networks. We define the dictionary

\begin{align}
    V^{(l)}_{acm} &= \{l: A_l \in  \mathbb{R}^{[|S|, n_l, 2l+1]} \}
\end{align}
Here $V^{(l)}_{acm}$ maps rotation order $l$ to a multidimensional array $A_l$ of inputs. Here we follow the following conventions on variables
\begin{enumerate}
\item $a$: The point index referring to a given atom
\item $c$: The channel index. Recall that a channel refers to a copy of an irreducible representation of order $l$. $n_l$ is used to refer to the number of channels.
\item $m$: The representation index, ranging from $1$ to $2l+1$, corresponding to the dimensions of the Wigner $D$-matrix
\end{enumerate}


\begin{figure}
    \centering
    % \includegraphics[width=0.7\textwidth]{figures/Differentiable Physics/Molecular ML/Tensor Field Networks/acm_features.png}
    \includegraphics[width=0.7\textwidth]{figures/Differentiable Physics/Molecular ML/Tensor Field Networks/Vacm dictionaries.png}
    \caption{The dictionaries $V^{(l)}_{acm}$ contain equivariant features used by the tensor field network and other related architectures.}
    \label{fig:acm_features}
\end{figure}

In the reminder of this chapter, we will explicitly construct some equivariant tensor field layers that operate on these dictionary objects.

\section{Point Convolutions}

We define point convolutions as objects that take as input dictionaries $V^{(l)}_{acm}$. These convolutions can be viewed as a generalization of a standard convolutional layer (with the usual convolution corresponding to the case where the atoms are arranged in a regular grid). 

\subsection{Spherical Filters}

Recall that the spherical harmonics are functions defined on the sphere

\begin{align}
    Y^{(l)}_m: S^2 \to \mathbb{R}
\end{align}
Here $l$ is again the rotation order. For order $0$ and order $1$, we have that the spherical harmonics are functions

\begin{align}
    Y^{(0)}(\hat{r}) &\propto 1 \\
    Y^{(1)}(\hat{r}) &\propto \hat{r} 
\end{align}

Spherical convolutions are equivariant to $SO(3)$ and for all $g \in SO(3)$, satisfy the equivariance relationship
    
\begin{align}
    Y^{(l)}_m(R(g)\hat{r}) &= \sum_{m'} D^{(l)}_{mm'}(g)Y^{(l)}_{m'}(\hat{r})
\end{align}

Here $D$ is the Wigner $D$-matrix. Spherical harmonic functions provide useful primitives to define rotation equivariant filters. Let $l_i$ denote the rotation order of the input and $l_f$ the rotation order of the filter. We mandate that filter $F^{(l_f, l_i)}$ takes the functional form

\begin{align}
    F^{(l_f, l_i)}_{cm}(\hat{r}) &= R_c^{(l_f, l_i)}(r) Y^{(l_f)}_m(\hat{r})
\end{align}
where $R_c^{(l_f, l_i)}$ depends only on the radius $r$ of $\hat{r}$ 
\begin{align}
R_c^{(l_f, l_i)}: \mathbb{R}^+ \to \mathbb{R}
\end{align} is a learnable function in this filter. This filter by construction is also equivariant to $SO(3)$. 

\subsection{Combining Representations with Tensor Products}

The tensor product provides a useful tool for combining representations. Suppose $\rho^X$ and $\rho^Y$ are two representations. Then $\rho^X \otimes \rho^Y$ is a representation on $X \otimes Y$. Suppose that that we have an irreducible representation $u^{(l_1)}$ of rotation order $l_1$ and an irreducible representation $v^{(l_2)}$ of rotation order $l_2$. Then $u^{(l_1)} \otimes v^{(l_2)}$ is a representation of order $l_1 \times l_2$. The coefficients of this combined representation are given explicitly by the Clebsch-Gordon coefficients
\begin{align}
(u \otimes v)^{(l)}_m &= \sum_{m_1=-l_1}^{l_1} \sum_{m_2=-l_2}^{l_2} C^{(l,m)}_{(l_1,m_1),(l_2,m_2)} u^{(l_1)}_{m_1} v^{(l_2)}_{m_2}
\end{align}

\subsection{Combining Clebsch-Gordon and Spherical Filters}
Given this definition, we define the point convolution layer. The learned spherical filters inhabits one representation, and the input inhabits another. To combine the filter representation and the input representation, we have to make use of the Clebsch-Gordon rule. Let $l_o$ denote the rotation order of the output. Then we define the point convolution layer by the rule

\begin{align}
L^{(l_0)}_{acm_o}(\vec{r}_a, V^{(l_i)}_{acm_i}) &= \sum_{m_f, m_i} C^{(l_o, m_o)}_{(l_f, m_f), (l_i, m_i)} \sum_{b \in S} F^{(l_f, l_i)}_{cm_f}(\vec{r}_{ab})V^{(l_i)}_{bcm_i}
\end{align}

The rule is dense to parse, but intuitively, it combines each filter across all pairs of atoms against the input to produce an output. This joint construction is equivariant to rotations.
\section{Self-Interaction Layers}

Self interaction layers are defined by the equation
\begin{align}
V^{(l)}_{acm} &= \sum_{c'} 
 W^{(l)}_{cc'}V^{(l)}_{ac'm}
\end{align}
These layer blend information across channels of $V^{(l)}_{acm}$. Because $W$ does not depend on the representation index $m$, it commutes with the Wigner $D$-matrices.

\section{Nonlinearities}
The nonlinear layers defined here act along the $m$ dimension. We define the norm operation
\begin{align}
\|V\|^{(l)}_{ac} &= \sqrt{\sum_m |V^{(l)}_{acm}|^2}
\end{align}
Suppose we have some learnable function for each $l$
\begin{align}
\eta^{(l)}: \mathbb{R} \to \mathbb{R}
\end{align}
We define the nonlinearity layer as
\begin{align}
V^{(l)}_{acm} &= \eta^{(l)}(\|V\|^{(l)}_{ac} + b^{(l)}_c) V^{(l)}_{acm}
\end{align}
where $b^{(l)}_c$ is a bias offset. The Wigner $D$-matrices are unitary, so they do not change the norm $\|V\|$, which means that our nonlinearities commute with the $D$-matrices.

\section{Building Tensor Field Architectures}

The point-convolution, self-interaction, and nonlinear layers can be combined and recombined as in standard convolutional architectures. Unlike standard convolutional architectures, features $V^{(l)}_{acm}$ hold information for different orders $l$ allowing for manipulation of scalar, vector, matrix, and higher order features.

\begin{figure}
    \centering
    % \includegraphics{figures/Differentiable Physics/Molecular ML/Tensor Field Networks/tfn_arch.png}
    \includegraphics{figures/Differentiable Physics/Molecular ML/Tensor Field Networks/Information flow_tf arch.png}
    \caption{An illustration of information flow in a 3-layer tensor field architecture with scalar and vector features \url{https://arxiv.org/pdf/1802.08219.pdf}}
    \label{fig:my_label}
\end{figure}

\section{Exercises}
\begin{enumerate}
\item Prove that the composition of two equivariant networks is equivariant.
\item Prove that if a network is equivariant with respect to transformation $h$ and transformation $g$ separately, it is equivariant with respect to composed transformation $hg$.
\item Prove that filter $F_c^{(l_f, l_i)}$ is equivariant to $SO(3)$.
\end{enumerate}

\end{document}